\chapter{Figures and Margins}\label{ch:figs}

\newthought{Tufte layouts offer} three figure sizes: margin, text width, and
full width. Pick by how much room the picture needs.

\section{A diagram in the margin (TikZ)}

\begin{marginfigure}
  \begin{tikzpicture}
    \node[box] (a) {Draft};
    \node[box,below=6mm of a] (b) {Edit};
    \node[box,below=6mm of b] (c) {Typeset};
    \draw[flow] (a) -- (b);
    \draw[flow] (b) -- (c);
    \draw[flow] (c.east) -- ++(8mm,0) |- (a.east);
  \end{tikzpicture}
  \caption{The writing loop.}\label{fig:loop}
\end{marginfigure}

Figure~\ref{fig:loop} sits in the margin beside the text that discusses it.
Every figure gets a label and a caption that makes sense on its own.

\section{A full-width chart (pgfplots, data from CSV)}

\begin{figure*}
  \begin{tikzpicture}
    \begin{axis}[book plot,width=\linewidth,height=5.5cm,
                 xlabel={Year},ylabel={Value},
                 xtick=data]
      \addplot table[col sep=comma,x=year,y=value]{data/sample.csv};
    \end{axis}
  \end{tikzpicture}
  \caption{A pgfplots chart that reads \texttt{data/sample.csv}. Edit the CSV
  and rebuild; nothing is copied by hand.}\label{fig:pgf}
\end{figure*}

\section{A generated figure (Python)}

For heavier charts, write a script in \texttt{figures/src/} and let the
Makefile render it to \texttt{figures/out/}.

\IfFileExists{figures/out/growth.pdf}{%
\begin{figure}
  \includegraphics[width=\linewidth]{growth}
  \caption{Rendered by \texttt{figures/src/growth.py} with the shared
  matplotlib style.}\label{fig:py}
\end{figure}%
}{\fbox{Run \texttt{make figures} to generate Figure~\ref{fig:py}.}}

\section{A table}

\begin{margintable}
  \footnotesize
  \begin{tabular}{ll}
    \toprule
    Mode & AI role \\
    \midrule
    Draft & off \\
    Edit & reviewer \\
    Typeset & agent \\
    \bottomrule
  \end{tabular}
  \caption{Working modes.}
\end{margintable}

Use \texttt{booktabs} rules only (no vertical lines), as in the margin table.
